\documentclass[instructions.tex]{subfiles}
\begin{document}
 
\chapter{De la patience dans l'adversité.}
 
On ne peut servir Dieu, on ne peut mériter le ciel, sans la patience; deux choses que nous apprend saint Paul, quand il dit : La patience vous est nécessaire, afin qu'en faisant la volonté de Dieu, vous remportiez l'effet de ses promesses\footnote{(1) Hebr. 10.}.
 
I. Pour servir Dieu, il faut être soumis à sa divine volonté. On n'est donc pas serviteur de Dieu, ni vraiment chrétien, quand on vit sans patience. Dieu veut que nous souffrions; et en nous appellant au christianisme, il nous appelle à la croix, parce qu'il nous appelle à la suite de Jésus crucifié. Si quelqu'un, dit Jésus-Christ, veut venir après moi, qu'il porte sa croix tous les jours, et qu'il me suive\footnote{(2) Matth. 16. 26,}.
 
Le vrai serviteur de Dieu n'est pas toujours celui qui fait de grandes choses, mais celui qui est disposé à souffrir. C'est beaucoup faire pour Dieu que de souffrir pour son amour. C'est donc peu de pratiquer les autres vertus, si l'on manque de patience; si toutefois on peut avoir quelques véritables vertus, quand on n'a pas celle-ci.
 
On voit des personnes fréquenter les sacremens, témoigner dans l'oraison de grands désirs d'aimer Dieu, former de grands desseins pour sa gloire; mais la moindre disgrace, une parole, un mépris, leur fait prendre feu et les irrite. Que cette vertu est fragile, qu'une parole fait échouer! L'on sert Dieu, l'on est en paix quand il console, quand on n'a rien à souffrir; mais, lorsqu'il nous fait part de sa croix, ce n'est plus la même chose. Que d'illusions dans ces prétendues vertus! Dieu n'est-il pas également saint et adorable, de quelque manière qu'il nous traite? Mérite-t-il moins notre soumission dans les disgraces que dans la prospérité?
 
Loin de là, rien ne doit plus nous attacher à Dieu que l'adversité; elle est le partage ordinaire des enfans de Dieu. Le Seigneur châtie ceux qu'il aime, dit saint Paul, et frappe tous ceux qu'il reçoit au nombre de ses enfans. Il dit tous et n'excepte personne, car y a-t-il un fils qui ne soit châtié de son père\footnote{(1) Quis enim filius quem non corripit pater? \textit{Hebr.} 12.}? Ainsi, continue saint Paul, si vous n'éprouvez pas les afflictions dont Dieu fait part à ses serviteurs, vous n'appartenez plus à Jésus-Christ, vous lui êtes étrangers, et non pas ses véritables enfans\footnote{(2) Si extra disciplinam estis, cujus participes facti sunt omnes, ergò adulteri, et non filii estis. \textit{Ibid.}}. Saint Ambroise eut donc raison de sortir d'une maison dont l'hôte, riche et scandaleux, n'avait jamais éprouvé aucune disgrâce. Sortons d'ici, dit le saint; cette maison périra, puisqu'elle est sans adversité. A peine fut-il sorti, que la maison s'écroula et écrasa l'hôte.
 
Les adversités ne sont cependant pas toujours une marque qu'on est du nombre des serviteurs de Dieu. Les méchans ont leurs contradictions et leurs disgraces; mais ils traînent leur croix plutôt qu'ils ne la portent; et, par le mauvais usage qu'ils en font, ils se perdent par le chemin qui devrait les sauver, à l'exemple du mauvais larron, qui de sa croix descendit aux enfers, tandis que le bon larron monta au ciel.
 
II. Si la patience est la marque des serviteurs de Dieu, elle est aussi le chemin qui conduit au ciel. Jésus-Christ nous a promis la gloire, mais il veut que nous l'ayons avec honneur, et que nous la gagnions par notre courage. Il ne nous a montré, pour aller au ciel, qu'un chemin hérissé d'épines; son Evangile ne nous prêche que la croix; il a voulu lui-même être crucifié, et n'est entré dans le ciel que par cette voie. Il a fallu, dit l'Evangile, que le Christ souffrît, afin qu'il entrât dans sa gloire\footnote{(1) Oportuit pati Christum, et ita intrare in gloriam suam. \textit{Luc.} 24.}.
 
Voilà ce que les apôtres prêchaient de sa part : C'est par beaucoup de tribulations qu'il faut que nous entrions dans le royaume de Dieu. Oui, il le faut; c'est une loi que le Fils de Dieu a subie. Serait-il juste que nous entrassions dans sa gloire par un chemin de roses, tandis qu'il n'y est entré lui-même que par un chemin d'épines\footnote{(2) Per multas tribulationes oportet nos intrare in regnum Dei. \textit{Act.} 14.}?
 
Ne nous y trompons pas : on ne va point de ce monde à la gloire du ciel par les délices. Ceux qui vivent dans les plaisirs et l'abondance, qui n'éprouvent ni adversités ni disgraces, semblent être réservés aux vengeances de Dieu, comme les victimes qu'on engraisse pour les sacrifier. C'est une vérité indubitable, que ceux qui ne veulent pas souffrir en ce monde souffriront en l'autre. Tout pécheur qui n'est pas purifié en cette vie par le feu de l'adversité ou de la pénitence, sera puni en l'autre par le feu de l'éternité.
 
Pourquoi donc craignez-vous de souffrir? Pourquoi tant d'impatience et de murmures? C'est que vous n'aimez que votre satisfaction, vos plaisirs et vos biens. Mais tous les plaisirs et tous les biens de la terre sont-ils comparables à un moment de délices du ciel, que vous vous exposez à perdre par vos impatiences? Ne vaut-il pas mieux souffrir quelques afflictions, quelques pertes pendant quelques jours, que d'être exclu du séjour de la gloire? Aimez-vous mieux être un jour précipité dans les feux, que de souffrir ici-bas quelque temps, pour mériter un bonheur sans fin?
 
L'homme le plus sensuel se priverait volontiers d'une liqueur agréable, s'il savait qu'il dût en être empoisonné, et rendrait grâces à celui qui la lui ôterait. Ah! insensé, vous savez que le trop grand amour de vous-même, l'attache à vos plaisirs et aux biens, empoisonnent votre âme, qu'ils embraseront l'enfer pour vous punir; et cependant vous vous livrez à l'emportement quand on vous en prive, ou quand on contredit vos inclinations. Où est votre raison? Si vous ne pouvez souffrir quelques momens, comment pourrez-vous souffrir toujours?
 
Les vrais serviteurs de Dieu, loin d'être en paix dans la prospérité, et de chercher les joies du monde, craignent pour leur salut dès qu'ils sont sans tribulations. Les souffrances et la croix font leur sûreté, ils les demandent comme des moyens d'arriver à la gloire. Ah! mon Dieu, disent-ils avec saint Augustin, affligez-moi, purifiez-moi, ne m'épargnez point en cette vie, pourvu que vous m'épargniez en l'autre, et que je vous possède en l'éternité\footnote{(1) Hic ure, hic seca, modò in æternum parcas. S. Aug.}.
 
\section{Résolutions}
 
\primo Persuadé que la patience m'est nécessaire presque à chaque instant de la vie, je la demanderai souvent à Dieu.
 
\secundo Et dès que je m'appercevrai qu'il s'élève en moi quelque mouvement contraire à cette vertu fragile, je m'efforcerai d'apaiser sur-le-champ ce mouvement.
 
\tertio A plus forte raison renoncerai-je aux murmures, aux sentimens de haine, aux désirs de vengeance, et autres excès de cette nature, aussitôt que ma conscience m'avertira qu'ils naissent dans mon cœur. 

\prayer{Mon Dieu, que de fautes j'ai faites jusqu'ici, parce que j'ai manqué de patience! Donnez-moi, je vous en supplie, cette vertu, qui est, en quelque sorte, le fondement de toutes les autres vertus.}
 
\end{document}
